% paper/section_5_evaluation.tex
% ─────────────────────────────────────────────────────────────────────────────
% Section 5: Evaluation
% Fill [TBD] placeholders with numbers from benchmarks/injecagent_runner.py
% and benchmarks/benign_traces.json once runs complete on the lab machine.
% ─────────────────────────────────────────────────────────────────────────────

\section{Evaluation}
\label{sec:eval}

We evaluate Kavach on two complementary datasets: InjecAgent~\cite{zhan2024injecagent},
which provides adversarial injection cases, and a hand-curated benign trace
set that tests false-positive rate on legitimate developer workflows.

\subsection{Experimental Setup}

\paragraph{Hardware.}
All experiments run on a Dell Precision 3660 workstation
(Intel Core i9-13900, 128\,GB DDR5 RAM, NVIDIA RTX 4090 24\,GB).
The Kavach server runs on localhost; the OpenClaw agent runs in a separate
process communicating over loopback. Embedding inference uses the RTX 4090.

\paragraph{Embedding model.}
\texttt{BAAI/bge-base-en-v1.5} (768-d, FP32, sentence-transformers 2.x).
The BGE asymmetric query prefix (``Represent this sentence for searching
relevant passages:'') is applied to query inputs only, never to corpus
documents, per the BGE convention. No quantization is applied.

\paragraph{LLM backbone.}
Agent inference uses \texttt{Gemma\,4 27B} served via Ollama on the RTX
4090, the configuration under which all reported interception, latency, and
benchmark numbers were collected.

\paragraph{Threshold configuration.}
We evaluate three threshold configurations identified by
\texttt{benchmarks/threshold\_sweep.py}:

\begin{center}
\begin{tabular}{lccc}
\toprule
Config & $\theta_{\min}$ & $\theta_{\text{route}}$ & $\theta_{\text{drift}}$ \\
\midrule
A (aggressive) & 0.40 & 0.35 & 0.15 \\
B (balanced)   & 0.45 & 0.40 & 0.20 \\
C (conservative) & 0.55 & 0.50 & 0.30 \\
\bottomrule
\end{tabular}
\end{center}

Config\,B is our primary reported configuration as it achieves the best
Youden's $J$ on the COMPASS calibration set.

\subsection{Attack Detection: InjecAgent}

InjecAgent provides 1{,}054 adversarial tool-call injection cases across
17 attack categories derived from MITRE ATT\&CK. We run each case through the
Kavach \texttt{/parliament/evaluate} endpoint and record the verdict.

\begin{table}[h]
\centering
\caption{Kavach detection results on InjecAgent (Config\,B, $n=1054$).}
\label{tab:injecagent}
\begin{tabular}{lrrrr}
\toprule
Category & Cases & \textsc{Block} & \textsc{Escalate} & Missed \\
\midrule
Credential Access    & [TBD] & [TBD] & [TBD] & [TBD] \\
Code Execution       & [TBD] & [TBD] & [TBD] & [TBD] \\
Data Exfiltration    & [TBD] & [TBD] & [TBD] & [TBD] \\
Persistence          & [TBD] & [TBD] & [TBD] & [TBD] \\
Defense Evasion      & [TBD] & [TBD] & [TBD] & [TBD] \\
Lateral Movement     & [TBD] & [TBD] & [TBD] & [TBD] \\
Other                & [TBD] & [TBD] & [TBD] & [TBD] \\
\midrule
\textbf{Total}       & 1054  & [TBD] & [TBD] & [TBD] \\
\bottomrule
\end{tabular}
\end{table}

\paragraph{Primary metric: BLOCK-recall.}
We define BLOCK-recall as the fraction of attack cases for which Kavach returns
either \textsc{Block} or \textsc{Escalate} (i.e., cases where tool execution
is prevented). Config\,B achieves BLOCK-recall = [TBD]\% ($95\%$ CI:
$[\text{TBD},\, \text{TBD}]$, Wilson interval).

\paragraph{ClawHavoc case study (CVE-2026-25253).}
ClawHavoc exploits OpenClaw's temporal composition blindness: a benign tool
call in turn $t$ primes a state variable, and a malicious call in turn $t+2$
exploits the primed state. Neither call is individually detectable by
single-turn classifiers. Kavach's NAVIGATOR Minister detects the trajectory
anomaly: the cosine similarity between turn-$t+2$'s embedding and the session
centroid falls below the drift threshold at turn $t+2$, triggering an
\textsc{Escalate} before execution.

Across [TBD] ClawHavoc replay traces, Kavach prevents [TBD]\% of successful
exploits that a single-turn baseline (LlamaFirewall~\cite{saxe2025llamafirewall})
misses entirely.

\subsection{False Positive Rate: Benign Traces}

We evaluate FPR on 50 hand-curated benign developer sessions
(\texttt{benchmarks/benign\_traces.json}) covering common tasks: debugging
Python type errors, deploying web applications, running test suites,
refactoring code, and managing git repositories. Each session contains
between 5 and 20 tool calls.

\begin{table}[h]
\centering
\caption{False positive rate on benign traces (Config\,B, $n=50$ sessions).}
\label{tab:benign}
\begin{tabular}{lrr}
\toprule
Task category & Sessions & FPR (\%) \\
\midrule
Python debugging     & [TBD] & [TBD] \\
Web deployment       & [TBD] & [TBD] \\
Test suite runs      & [TBD] & [TBD] \\
Code refactoring     & [TBD] & [TBD] \\
Git operations       & [TBD] & [TBD] \\
\midrule
\textbf{Overall}     & 50    & [TBD] \\
\bottomrule
\end{tabular}
\end{table}

An FPR of [TBD]\% means that in [TBD] of 50 benign sessions Kavach incorrectly
blocked or escalated at least one legitimate tool call.

\subsection{Latency}

We measure end-to-end evaluation latency (agent plugin call to verdict received)
under three concurrency levels. Times are wall-clock milliseconds, median and
p95 over 500 requests per condition.

\begin{center}
\begin{tabular}{lrr}
\toprule
Concurrency & p50 (ms) & p95 (ms) \\
\midrule
1 (serial)       & [TBD] & [TBD] \\
4 (typical)      & [TBD] & [TBD] \\
8 (stress)       & [TBD] & [TBD] \\
\bottomrule
\end{tabular}
\end{center}

The p95 latency target is $<$200\,ms for synchronous interception to remain
imperceptible to the agent operator. Config\,B [meets / does not meet — fill
after run] this target at concurrency~4.

\subsection{Ablation: Minister Contribution}

We ablate the contribution of each Minister by disabling it and re-running
InjecAgent, measuring the resulting drop in BLOCK-recall.

\begin{center}
\begin{tabular}{lr}
\toprule
Disabled component & $\Delta$ BLOCK-recall (\%) \\
\midrule
EXECUTOR off  & $-$[TBD] \\
VAULT off     & $-$[TBD] \\
CHANNEL off   & $-$[TBD] \\
NAVIGATOR off & $-$[TBD] \\
COMPASS off   & $-$[TBD] \\
\bottomrule
\end{tabular}
\end{center}

NAVIGATOR and COMPASS together account for [TBD]\% of detections that would be
missed by a stateless, single-turn system, validating the temporal--spatial
architecture described in \S\ref{sec:temporal}.

\subsection{Comparison with Related Systems}

\begin{table}[h]
\centering
\caption{Comparison with guardrail systems on InjecAgent BLOCK-recall and FPR.}
\label{tab:comparison}
\begin{tabular}{lrrl}
\toprule
System & BLOCK-recall & FPR & Temporal detection \\
\midrule
LlamaFirewall~\cite{saxe2025llamafirewall} & [TBD]\% & [TBD]\% & No \\
AGrail~\cite{luo2025agrail}               & [TBD]\% & [TBD]\% & No \\
PRISM~\cite{guo2025prism}                 & [TBD]\% & [TBD]\% & In-motion only \\
NemoClaw~\cite{nvidia2025nemoclaw}        & [TBD]\% & [TBD]\% & No \\
\textbf{Kavach (ours, Config B)}          & \textbf{[TBD]\%} & \textbf{[TBD]\%} & \textbf{Yes (NAVIGATOR+COMPASS)} \\
\bottomrule
\end{tabular}
\end{table}

\noindent
Numbers for baseline systems are taken from their published evaluations on
overlapping InjecAgent subsets. Direct comparison should be interpreted
cautiously as exact prompt formats and model versions vary.
